\section*{Capítulo \ref{ch:g9:acids-bases-ph} --- Ácidos, bases y pH}

\begin{solution}{exo:g9:acids-bases-ph:1}
\omterm{def:g9:acids-bases-ph:ph}{pH} 3: ácida. \omterm{def:g9:acids-bases-ph:ph}{pH} 7: neutra. \omterm{def:g9:acids-bases-ph:ph}{pH} 11: básica.
\end{solution}

\begin{solution}{exo:g9:acids-bases-ph:2}
En una ácida, los \omterm{def:g9:ions:ion}{iones} \ce{H+}. En una básica, los \omterm{def:g9:ions:ion}{iones} \ce{OH-}.
\end{solution}

\begin{solution}{exo:g9:acids-bases-ph:3}
El zumo de lima (2), el café (5), el agua pura (7), el agua de mar
(8.1), el amoniaco doméstico (12).
\end{solution}

\begin{solution}{exo:g9:acids-bases-ph:4}
Se toca una tira, colocada en un recipiente limpio y seco, con una gota
de la disolución llevada con una varilla de vidrio limpia; se compara
enseguida el color con la escala.
\end{solution}

\begin{solution}{exo:g9:acids-bases-ph:5}
Alrededor de 5 tras una dilución por diez; alrededor de 6 tras una
dilución por cien.
\end{solution}

\begin{solution}{exo:g9:acids-bases-ph:6}
No. Diluir acerca el \omterm{def:g9:acids-bases-ph:ph}{pH} a 7, pero nunca lo hace pasar de 7: el agua es
neutra, y añadirla no puede hacer que los \omterm{def:g9:ions:ion}{iones} \ce{OH-} superen en
número a los \omterm{def:g9:ions:ion}{iones} \ce{H+}.
\end{solution}

\begin{solution}{exo:g9:acids-bases-ph:7}
La \omterm{def:g6:pure-substances-and-mixtures:mixture}{mezcla} desprende calor. El agua vertida sobre un ácido concentrado
podría hervir de golpe y lanzar ácido fuera; el ácido vertido despacio
sobre el agua se reparte enseguida en un volumen grande.
\end{solution}

\begin{solution}{exo:g9:acids-bases-ph:8}
Corrosivo: quema la piel y los ojos, ataca los metales. Llevar gafas y
guantes; enjuagar cualquier salpicadura con mucha agua.
\end{solution}

\begin{solution}{exo:g9:acids-bases-ph:9}
Tres diluciones (de 1 a 2, de 2 a 3, de 3 a 4): $10 \times 10 \times 10 =
1000$ veces.
\end{solution}

\begin{solution}{exo:g9:acids-bases-ph:10}
La leche de magnesia es básica (\omterm{def:g9:acids-bases-ph:ph}{pH} de alrededor de 10.5): reacciona con
una parte del ácido en exceso del estómago.
\end{solution}

\begin{solution}{exo:g9:acids-bases-ph:11}
No: con un \omterm{def:g9:acids-bases-ph:ph}{pH} de 8.1 sigue siendo ligeramente básico. Evoluciona hacia
un \omterm{def:g9:acids-bases-ph:ph}{pH} más bajo, menos básico, porque el dióxido de carbono del \omterm{def:g4:air-a-mixture-of-gases:air}{aire} se
\omterm{def:g2:mixing-and-dissolving:dissolve}{disuelve} en él.
\end{solution}

\begin{solution}{exo:g9:acids-bases-ph:12}
El \omterm{def:g9:acids-bases-ph:ph-paper}{papel indicador} se lee a simple vista comparándolo con una escala de
colores, con una precisión de alrededor de una unidad; el \omterm{def:g9:acids-bases-ph:ph-paper}{pHmetro} da un
número con uno o dos decimales. El \omterm{def:g9:acids-bases-ph:ph-paper}{pHmetro} es más preciso; los dos
resultados concuerdan dentro de la precisión del papel.
\end{solution}

\begin{solution}{pb:g9:acids-bases-ph:1}
\textbf{1.} Ácida: contiene sobre todo \omterm{def:g9:ions:ion}{iones} \ce{H+(aq)} y
\ce{Cl-(aq)}.

\textbf{2.} GHS05, corrosivo: un ácido ataca la piel, los ojos y los
metales.

\textbf{3.} Da un número preciso en lugar de un color que hay que
comparar a simple vista.

\textbf{4.} Diez veces más: una unidad de \omterm{def:g9:acids-bases-ph:ph}{pH} corresponde a una dilución
por diez.

\textbf{5.} \qty{10}{L}, de \omterm{def:g9:acids-bases-ph:ph}{pH} 3.

\textbf{6.} Tres: de \omterm{def:g9:acids-bases-ph:ph}{pH} 2 a 3, de 3 a 4 y de 4 a 5.

\textbf{7.} $1 \times 10 \times 10 \times 10 = \qty{1000}{L}$.

\textbf{8.} No: diluir solo acerca el \omterm{def:g9:acids-bases-ph:ph}{pH} a 7.

\textbf{9.} \ce{H+ + OH- -> H2O}.

\textbf{10.} El ácido se destruye, se convierte en agua, en lugar de
repartirse en un volumen enorme; basta un poco de base, y no hacen falta
mil litros de agua.

\textbf{11.} Bicarbonato disuelto en agua (o agua jabonosa).

\textbf{12.} $1000 - 1 = \qty{999}{L}$ de agua.
\end{solution}
